% Options for packages loaded elsewhere
\PassOptionsToPackage{unicode}{hyperref}
\PassOptionsToPackage{hyphens}{url}
\PassOptionsToPackage{dvipsnames,svgnames,x11names}{xcolor}
\documentclass[
  11pt,
]{article}
\usepackage{xcolor}
\usepackage[a4paper,margin=19mm]{geometry}
\usepackage{amsmath,amssymb}
\setcounter{secnumdepth}{-\maxdimen} % remove section numbering
\usepackage{iftex}
\ifPDFTeX
  \usepackage[T1]{fontenc}
  \usepackage[utf8]{inputenc}
  \usepackage{textcomp} % provide euro and other symbols
\else % if luatex or xetex
  \usepackage{unicode-math} % this also loads fontspec
  \defaultfontfeatures{Scale=MatchLowercase}
  \defaultfontfeatures[\rmfamily]{Ligatures=TeX,Scale=1}
\fi
\usepackage{lmodern}
\ifPDFTeX\else
  % xetex/luatex font selection
  \setmainfont[]{Be Vietnam Pro}
  \setmonofont[]{Menlo}
\fi
% Use upquote if available, for straight quotes in verbatim environments
\IfFileExists{upquote.sty}{\usepackage{upquote}}{}
\IfFileExists{microtype.sty}{% use microtype if available
  \usepackage[]{microtype}
  \UseMicrotypeSet[protrusion]{basicmath} % disable protrusion for tt fonts
}{}
\makeatletter
\@ifundefined{KOMAClassName}{% if non-KOMA class
  \IfFileExists{parskip.sty}{%
    \usepackage{parskip}
  }{% else
    \setlength{\parindent}{0pt}
    \setlength{\parskip}{6pt plus 2pt minus 1pt}}
}{% if KOMA class
  \KOMAoptions{parskip=half}}
\makeatother
\setlength{\emergencystretch}{3em} % prevent overfull lines
\providecommand{\tightlist}{%
  \setlength{\itemsep}{0pt}\setlength{\parskip}{0pt}}

\usepackage{newunicodechar}
\newunicodechar{→}{\ensuremath{\rightarrow}}
\newunicodechar{↔}{\ensuremath{\leftrightarrow}}
\usepackage{fvextra}
\DefineVerbatimEnvironment{verbatim}{Verbatim}{breaklines=true,breakanywhere=true,fontsize=\small}
\usepackage{fancyhdr}
\usepackage{array}
\pagestyle{fancy}\fancyhf{}
\fancyhead[L]{\small MovieLOD · Bài làm hoàn chỉnh}
\fancyfoot[C]{\small\thepage}\setlength{\headheight}{15pt}
\renewcommand{\arraystretch}{1.14}\setlength{\extrarowheight}{1pt}
\AtBeginEnvironment{longtable}{\small}
\setlength{\emergencystretch}{3em}\setlength{\parskip}{4pt}
\usepackage{titling}\setlength{\droptitle}{-1.5cm}
\pretitle{\begin{center}\Large\bfseries\color{teal}}
\posttitle{\par\end{center}\vspace{-0.5em}}
\predate{\begin{center}\small}\postdate{\par\end{center}\vspace{-1em}}
\usepackage{bookmark}
\IfFileExists{xurl.sty}{\usepackage{xurl}}{} % add URL line breaks if available
\urlstyle{same}
% fallback for those not using the hyperref driver hyperxmp:
\makeatletter
\@ifundefined{xmpquote}{\newcommand{\xmpquote}[1]{#1}}{}
\makeatother
\hypersetup{
  pdftitle={MovieLOD: checklist ảnh minh chứng Protégé},
  colorlinks=true,
  linkcolor={teal},
  filecolor={Maroon},
  citecolor={Blue},
  urlcolor={teal},
  pdfcreator={LaTeX via pandoc}}

\title{MovieLOD: checklist ảnh minh chứng Protégé}
\author{}
\date{11 vị trí trong bộ 24 slide · 09/10/2026}

\begin{document}
\maketitle

\subsection{Cách chèn}\label{cuxe1ch-chuxe8n}

Có 11 khung ảnh ghi rõ mã, tên file và thao tác ngay trên slide. Lần tạo
này không có ảnh Protégé thật: chưa có ảnh giao diện đã xác minh cho
ontology mới; không dùng ảnh dựng làm screenshot. Ảnh ứng dụng Web trong
slide là ảnh chụp thực.

Cách 1: chèn ảnh vào PowerPoint và che/xóa khung chờ cùng phần hướng dẫn
của khung đó. Cách 2: lưu đúng tên ở
\texttt{evidence\allowbreak{}/\allowbreak{}protege\allowbreak{}/\allowbreak{}}
(\allowbreak{}P\allowbreak{}N\allowbreak{}G/\allowbreak{}\allowbreak{}J\allowbreak{}P\allowbreak{}G\allowbreak{}/\allowbreak{}\allowbreak{}J\allowbreak{}P\allowbreak{}E\allowbreak{}G
cùng stem), chạy
\texttt{.\allowbreak{}venv/\allowbreak{}bi\allowbreak{}n/\allowbreak{}python\allowbreak{} src/\allowbreak{}mak\allowbreak{}e\_\allowbreak{}slides\allowbreak{}\_\allowbreak{}video.\allowbreak{}p\allowbreak{}y -\allowbreak{}-\allowbreak{}slid\allowbreak{}es-\allowbreak{}only};
khung chờ tự thay bằng ảnh. Speaker Notes giữ quy trình chụp.

Chụp đúng cửa sổ/view, chữ đủ lớn (gợi ý 1440×900 trở lên); không lấy cả
desktop có ứng dụng khác. Giữ tên ontology và IRI/thông tin cần chứng
minh. Ảnh cây khai báo không được ghi nhãn inferred.

\subsection{Danh sách ảnh}\label{danh-suxe1ch-ux1ea3nh}

\subsubsection{P00 --- Slide 05: IRI và phiên bản
3.0.0}\label{p00-slide-05-iri-vuxe0-phiuxean-bux1ea3n-3.0.0}

\textbf{Tên file:}
\texttt{P00\_\allowbreak{}onto\allowbreak{}logy\_\allowbreak{}hea\allowbreak{}der.\allowbreak{}png}.
\textbf{Mở:}
\texttt{ontology\allowbreak{}/\allowbreak{}\allowbreak{}Movie\_\allowbreak{}\allowbreak{}K\allowbreak{}nowledge\allowbreak{}\_\allowbreak{}\allowbreak{}Graph.\allowbreak{}o\allowbreak{}wl}.

\begin{enumerate}
\def\labelenumi{\arabic{enumi}.}
\tightlist
\item
  Mở
  ontology\allowbreak{}/\allowbreak{}\allowbreak{}Movie\_\allowbreak{}\allowbreak{}K\allowbreak{}nowledge\allowbreak{}\_\allowbreak{}\allowbreak{}Graph.\allowbreak{}o\allowbreak{}wl.\allowbreak{}
\item
  Chọn Active Ontology / Ontology Header.
\item
  Giữ IRI và versionInfo 3.0.0 trong ảnh.
\end{enumerate}

\textbf{Mục cần thấy:} 37 lớp có tên; không đếm biểu thức anonymous như
lớp có tên.

\subsubsection{P01 --- Slide 06: Cây Work /
Film}\label{p01-slide-06-cuxe2y-work-film}

\textbf{Tên file:}
\texttt{P01\_\allowbreak{}film\allowbreak{}\_\allowbreak{}hierarc\allowbreak{}hy.\allowbreak{}png}.
\textbf{Mở:}
\texttt{ontology\allowbreak{}/\allowbreak{}\allowbreak{}Movie\_\allowbreak{}\allowbreak{}K\allowbreak{}nowledge\allowbreak{}\_\allowbreak{}\allowbreak{}Graph.\allowbreak{}o\allowbreak{}wl}.

\begin{enumerate}
\def\labelenumi{\arabic{enumi}.}
\tightlist
\item
  Mở Classes → dbo:Work → dbo:Film.
\item
  Hiện ActionFilm,
  Award\allowbreak{}Win\allowbreak{}ning\allowbreak{}Film\allowbreak{}
  và các lớp giao.
\item
  Giữ IRI DBpedia của Film trong Description.
\end{enumerate}

\textbf{Mục cần thấy:} Không có
Feature\allowbreak{}F\allowbreak{}ilm/\allowbreak{}\allowbreak{}Anim\allowbreak{}ated\allowbreak{}Film\allowbreak{}/\allowbreak{}\allowbreak{}Documen\allowbreak{}tary\allowbreak{}Film\allowbreak{}
trong module cuối.

\subsubsection{P02 --- Slide 07: Cây chủ thể
DBpedia}\label{p02-slide-07-cuxe2y-chux1ee7-thux1ec3-dbpedia}

\textbf{Tên file:}
\texttt{P02\_\allowbreak{}agen\allowbreak{}t\_\allowbreak{}hierar\allowbreak{}chy.\allowbreak{}png}.
\textbf{Mở:}
\texttt{ontology\allowbreak{}/\allowbreak{}\allowbreak{}Movie\_\allowbreak{}\allowbreak{}K\allowbreak{}nowledge\allowbreak{}\_\allowbreak{}\allowbreak{}Graph.\allowbreak{}o\allowbreak{}wl}.

\begin{enumerate}
\def\labelenumi{\arabic{enumi}.}
\tightlist
\item
  Mở dbo:Agent → dbo:Person và
  dbo:\allowbreak{}\allowbreak{}Orga\allowbreak{}nisation\allowbreak{}.\allowbreak{}
\item
  Hiện Artist → Actor; Organisation → Company.
\item
  Mở Movie\allowbreak{}Dir\allowbreak{}ector / Writer →
  Screen\allowbreak{}Wr\allowbreak{}iter.\allowbreak{}
\end{enumerate}

\textbf{Mục cần thấy:} Reuse lớp DBpedia; không tạo ex:Actor hoặc
ex:\allowbreak{}\allowbreak{}Organ\allowbreak{}ization.\allowbreak{}

\subsubsection{P03 --- Slide 08: MovieGenre và
Award}\label{p03-slide-08-moviegenre-vuxe0-award}

\textbf{Tên file:}
\texttt{P03\_\allowbreak{}genr\allowbreak{}e\_\allowbreak{}award\_\allowbreak{}hierarch\allowbreak{}y.\allowbreak{}png}.
\textbf{Mở:}
\texttt{ontology\allowbreak{}/\allowbreak{}\allowbreak{}Movie\_\allowbreak{}\allowbreak{}K\allowbreak{}nowledge\allowbreak{}\_\allowbreak{}\allowbreak{}Graph.\allowbreak{}o\allowbreak{}wl}.

\begin{enumerate}
\def\labelenumi{\arabic{enumi}.}
\tightlist
\item
  Mở dbo:Genre →
  dbo:\allowbreak{}\allowbreak{}Movi\allowbreak{}e\allowbreak{}Genre.\allowbreak{}
\item
  Hiện ActionGenre / DramaGenre; chọn dbo:Award.
\item
  Giữ Description để đọc IRI và lớp cha.
\end{enumerate}

\textbf{Mục cần thấy:} Không có FictionGenre hoặc các nhóm award suy
đoán theo nhãn.

\subsubsection{P04 --- Slide 09: Ràng buộc
Contribution}\label{p04-slide-09-ruxe0ng-buux1ed9c-contribution}

\textbf{Tên file:}
\texttt{P04\_\allowbreak{}cont\allowbreak{}ribution\allowbreak{}\_\allowbreak{}restric\allowbreak{}tions.\allowbreak{}pn\allowbreak{}g}.
\textbf{Mở:}
\texttt{ontology\allowbreak{}/\allowbreak{}\allowbreak{}Movie\_\allowbreak{}\allowbreak{}K\allowbreak{}nowledge\allowbreak{}\_\allowbreak{}\allowbreak{}Graph.\allowbreak{}o\allowbreak{}wl}.

\begin{enumerate}
\def\labelenumi{\arabic{enumi}.}
\tightlist
\item
  Chọn
  ex:\allowbreak{}\allowbreak{}Contr\allowbreak{}ibution.\allowbreak{}
\item
  Hiện exactly 1 cho contribu\allowbreak{}tion\allowbreak{}By,
  contribu\allowbreak{}tion\allowbreak{}To, hasRole.
\item
  Chọn hasRole để thấy Functional và role individual.
\end{enumerate}

\textbf{Mục cần thấy:} Ba endpoint đúng 1; dữ liệu thiếu cần structural
validation riêng.

\subsubsection{P05 --- Slide 10: Domain / range /
inverse}\label{p05-slide-10-domain-range-inverse}

\textbf{Tên file:}
\texttt{P05\_\allowbreak{}obje\allowbreak{}ct\_\allowbreak{}prope\allowbreak{}rty.\allowbreak{}png}.
\textbf{Mở:}
\texttt{ontology\allowbreak{}/\allowbreak{}\allowbreak{}Movie\_\allowbreak{}\allowbreak{}K\allowbreak{}nowledge\allowbreak{}\_\allowbreak{}\allowbreak{}Graph.\allowbreak{}o\allowbreak{}wl}.

\begin{enumerate}
\def\labelenumi{\arabic{enumi}.}
\tightlist
\item
  Chọn contribu\allowbreak{}tion\allowbreak{}By:\allowbreak{}
  Contribution → dbo:Person.
\item
  Hiện Functional và inverse
  has\allowbreak{}Contr\allowbreak{}ibution.\allowbreak{}
\item
  Có thể chụp contribu\allowbreak{}ted\allowbreak{}To với property
  chain.
\end{enumerate}

\textbf{Mục cần thấy:} Chain has\allowbreak{}Contr\allowbreak{}ibution
rồi contribu\allowbreak{}tion\allowbreak{}To; không đặt cardinality trên
contribu\allowbreak{}ted\allowbreak{}To.\allowbreak{}

\subsubsection{P06 --- Slide 13: Định nghĩa
Filmmaker}\label{p06-slide-13-ux111ux1ecbnh-nghux129a-filmmaker}

\textbf{Tên file:}
\texttt{P06\_\allowbreak{}film\allowbreak{}maker\_\allowbreak{}eq\allowbreak{}uivalent\allowbreak{}\_\allowbreak{}class.\allowbreak{}p\allowbreak{}ng}.
\textbf{Mở:}
\texttt{ontology\allowbreak{}/\allowbreak{}\allowbreak{}Movie\_\allowbreak{}\allowbreak{}K\allowbreak{}nowledge\allowbreak{}\_\allowbreak{}\allowbreak{}Graph.\allowbreak{}o\allowbreak{}wl}.

\begin{enumerate}
\def\labelenumi{\arabic{enumi}.}
\tightlist
\item
  Chọn ex:\allowbreak{}\allowbreak{}Filmm\allowbreak{}aker.\allowbreak{}
\item
  Hiện Equivalent To với Person AND các nhánh SOME.
\item
  Đọc các nhánh Directing / Writing /
  Producin\allowbreak{}g\allowbreak{}Contrib\allowbreak{}ution.\allowbreak{}
\end{enumerate}

\textbf{Mục cần thấy:} Công thức không tự chứng minh reasoner đã chạy.

\subsubsection{P07 --- Slide 14: Cardinality trên
Contribution}\label{p07-slide-14-cardinality-truxean-contribution}

\textbf{Tên file:}
\texttt{P07\_\allowbreak{}thre\allowbreak{}e\_\allowbreak{}credit\allowbreak{}\_\allowbreak{}cardina\allowbreak{}lity.\allowbreak{}png\allowbreak{}}.
\textbf{Mở:}
\texttt{ontology\allowbreak{}/\allowbreak{}\allowbreak{}Movie\_\allowbreak{}\allowbreak{}K\allowbreak{}nowledge\allowbreak{}\_\allowbreak{}\allowbreak{}Graph.\allowbreak{}o\allowbreak{}wl}.

\begin{enumerate}
\def\labelenumi{\arabic{enumi}.}
\tightlist
\item
  Chọn
  ex:\allowbreak{}\allowbreak{}Three\allowbreak{}\allowbreak{}Credit\allowbreak{}Co\allowbreak{}ntributo\allowbreak{}r;
  hiện min 3 has\allowbreak{}Contr\allowbreak{}ibution
  Contribu\allowbreak{}tion.\allowbreak{}
\item
  Chạy HermiT; chọn Nolan ở inferred view.
\item
  Chụp thêm Functional hasRole và AllDifferent của bốn role.
\end{enumerate}

\textbf{Mục cần thấy:} 7 người đạt min 3; 17 người đạt min 2; không phải
COUNT DISTINCT.

\subsubsection{P08 --- Slide 16: Thời lượng theo
DBpedia}\label{p08-slide-16-thux1eddi-lux1b0ux1ee3ng-theo-dbpedia}

\textbf{Tên file:}
\texttt{P08\_\allowbreak{}runt\allowbreak{}ime\_\allowbreak{}data\allowbreak{}type.\allowbreak{}png\allowbreak{}}.
\textbf{Mở:}
\texttt{ontology\allowbreak{}/\allowbreak{}\allowbreak{}Movie\_\allowbreak{}\allowbreak{}K\allowbreak{}nowledge\allowbreak{}\_\allowbreak{}\allowbreak{}Graph.\allowbreak{}o\allowbreak{}wl}.

\begin{enumerate}
\def\labelenumi{\arabic{enumi}.}
\tightlist
\item
  Chọn dbo:runtime trong Data properties.
\item
  Hiện domain dbo:Work và range xsd:double.
\item
  Chọn Inception: runtime = 8880 giây.
\end{enumerate}

\textbf{Mục cần thấy:} 8880 giây = 148 phút; rdfs:label là annotation
property.

\subsubsection{P09 --- Slide 18: Cá thể Inception bản
mới}\label{p09-slide-18-cuxe1-thux1ec3-inception-bux1ea3n-mux1edbi}

\textbf{Tên file:}
\texttt{P09\_\allowbreak{}ince\allowbreak{}ption\_\allowbreak{}in\allowbreak{}dividual\allowbreak{}.\allowbreak{}png}.
\textbf{Mở:}
\texttt{ontology\allowbreak{}/\allowbreak{}\allowbreak{}Movie\_\allowbreak{}\allowbreak{}K\allowbreak{}nowledge\allowbreak{}\_\allowbreak{}\allowbreak{}Graph.\allowbreak{}o\allowbreak{}wl}.

\begin{enumerate}
\def\labelenumi{\arabic{enumi}.}
\tightlist
\item
  Mở
  Movie\_\allowbreak{}\allowbreak{}Kn\allowbreak{}owledge\_\allowbreak{}\allowbreak{}Graph.\allowbreak{}ow\allowbreak{}l
  3.0.0.
\item
  Chọn film-Q25188; hiện label, releaseYear 2010, runtime 8880.
\item
  Hiện dbo:\allowbreak{}dire\allowbreak{}ctor, genre,
  producti\allowbreak{}on\allowbreak{}Compan\allowbreak{}y và award.
\end{enumerate}

\textbf{Mục cần thấy:} Cá thể trong OWL khớp trang resource 3.0.0;
runtime là 8880 giây.

\subsubsection{P10 --- Slide 21: Các type suy luận của
Nolan}\label{p10-slide-21-cuxe1c-type-suy-luux1eadn-cux1ee7a-nolan}

\textbf{Tên file:}
\texttt{P10\_\allowbreak{}nola\allowbreak{}n\_\allowbreak{}inferr\allowbreak{}ed\_\allowbreak{}types\allowbreak{}.\allowbreak{}png}.
\textbf{Mở:}
\texttt{ontology\allowbreak{}/\allowbreak{}\allowbreak{}Movie\_\allowbreak{}\allowbreak{}K\allowbreak{}nowledge\allowbreak{}\_\allowbreak{}\allowbreak{}Graph.\allowbreak{}o\allowbreak{}wl}.

\begin{enumerate}
\def\labelenumi{\arabic{enumi}.}
\tightlist
\item
  Chọn HermiT → Start reasoner, chờ hoàn tất.
\item
  Chọn person-\allowbreak{}\allowbreak{}Q\allowbreak{}25191; mở inferred
  types.
\item
  Giữ Writer\allowbreak{}Di\allowbreak{}rector, Filmmaker,
  Three\allowbreak{}Cre\allowbreak{}dit\allowbreak{}Contr\allowbreak{}ibutor
  trong ảnh.
\end{enumerate}

\textbf{Mục cần thấy:} Type mới không được gán trực tiếp trong asserted
graph; giữ trạng thái reasoner.

\subsection{Kiểm tra trước khi dùng làm minh
chứng}\label{kiux1ec3m-tra-trux1b0ux1edbc-khi-duxf9ng-luxe0m-minh-chux1ee9ng}

\begin{itemize}
\tightlist
\item
  Đúng ontology 3.0.0 của bài; không dùng ảnh 2.0 làm minh chứng mô hình
  mới.
\item
  Ảnh cá thể dùng Knowledge Graph; ảnh mô hình dùng schema OWL.
\item
  IRI lớp DBpedia, cardinality và datatype đọc được.
\item
  Không nói
  Hermi\allowbreak{}T/\allowbreak{}\allowbreak{}P\allowbreak{}ellet đã
  chứng minh DL nếu chỉ chụp công thức hoặc chọn menu.
\item
  Nếu chạy reasoner, giữ tên và trạng thái; ghi cả lỗi nếu có.
\end{itemize}

Nguồn hướng dẫn giao diện:
\href{https://protegeproject.github.io/protege/views/}{Protégé Views}.
Ngữ nghĩa cần nhớ: \href{https://www.w3.org/TR/owl2-primer/}{OWL 2
Primer}. Nội dung lớp/thuộc tính/cá thể lấy từ OWL của bài.

\end{document}
